\documentclass[11pt,a4paper]{article}
\usepackage[UTF8,fontset=fandol]{ctex}
\usepackage{amsmath,amssymb,amsthm,mathtools,geometry,booktabs,array,hyperref}
\geometry{margin=24mm}
\hypersetup{colorlinks=true,linkcolor=blue,citecolor=blue,urlcolor=blue}
\hypersetup{pdftitle={A proof of Sun's ninety-dollar series for 1/pi},pdfauthor={李想}}
\setlength{\emergencystretch}{2em}
\renewcommand{\abstractname}{Abstract}
\renewcommand{\refname}{References}
\newtheorem{theorem}{Theorem}
\newtheorem{lemma}[theorem]{Lemma}
\newtheorem{proposition}[theorem]{Proposition}
\theoremstyle{remark}
\newtheorem{remark}[theorem]{Remark}
\newcommand{\tht}{\vartheta}
\newcommand{\dd}{\mathrm d}
\newcommand{\HH}{\mathcal H}
\newcommand{\code}[1]{\path{#1}}
\DeclareMathOperator{\Sym}{Sym}
\title{A proof of Sun's ninety-dollar series for $1/\pi$}
\author{李想}
\date{16 October 2026}

\begin{document}
\maketitle
\begin{center}
\small\textbf{Version 2 --- public preprint; independent review pending}
\end{center}
\begin{abstract}
We prove the double-series identity for $90/\pi$ proposed by Zhi-Wei Sun
in 2011. A restriction of its bivariate generating function satisfies a
third-order differential equation and an additional first-order gradient
identity. An algebraic substitution identifies the restricted period with
a modular form of level $60$. We prove the required general modular
parametrization independently from eta products, cusp orders, and exact
Sturm certificates of bounds $48$ and $96$. We also certify the degree-seven
modular equation by Sturm's theorem, isolate the complex multiplication
value by rational inequalities, and derive the weighted period evaluation
directly by differentiating the Fricke transformation. No tabulated
modular parametrization or special value is assumed. Exact rational-function
and integer-series certificates accompany the proof.
\end{abstract}

\section{The identity and the proof strategy}
Sun's conjectural identity \cite[Conjecture 5, (5.9)]{Sun} is the following.
\begin{theorem}\label{thm:main}
The series below converges absolutely, and
\begin{equation}\label{eq:main}
 \sum_{n=0}^{\infty}\frac{357n+103}{2160^n}\binom{2n}{n}
 \sum_{k=0}^{n}\binom nk\binom{n+2k}{2k}\binom{2k}{k}
                 (-324)^{n-k}=\frac{90}{\pi}.
\end{equation}
\end{theorem}
Sun offered ninety US dollars for the first rigorous proof. The identity
was also listed as an open conjecture in the 2024 version of
\cite[Section 4, (4.3)]{vHTY}. This article supplies a proof and its
computational certificates; it makes no claim that priority or a prize
payment has been confirmed by the sponsor.

For any variable $x$, write $\tht_x=x\dd/\dd x$.
The eta quotients and polynomials suggested by
\cite[Tables 2--4]{CCY} provided candidate modular data; no identity
from those tables is a theorem input here. Section~\ref{sec:modparam}
proves their required identities independently. The standard modular-form
inputs are the eta-quotient modularity and cusp-order criteria
\cite{Newman,Ligozat}, the eta transformation under $\tau\mapsto-1/\tau$,
and Sturm's theorem \cite{Sturm}. The needed derivative covariance is
proved explicitly below. The CM value and weighted evaluation are also
proved independently. The algebraic bridge consists of a restriction
identity, a gradient relation, and an exact transformation of differential
operators, each checked over a rational-function field.

\section{A bivariate period and its restriction}
Define, initially near $(u,v)=(0,0)$,
\begin{equation}\label{eq:H}
 \HH(u,v)=\sum_{j,k\geq0}
 \binom{2j+2k}{j+k}\frac{(j+3k)!}{j!(k!)^3}u^jv^k.
\end{equation}
Put $U=\tht_u\HH$ and $V=\tht_v\HH$. The substitution $j=n-k$
shows that the left side of \eqref{eq:main} is
\begin{equation}\label{eq:targetH}
 \bigl[357(U+V)+103\HH\bigr]
             \left(-\frac3{20},\frac1{2160}\right).
\end{equation}
Our restriction parameter is $q$, which is distinct from the modular
nome used later. Set
\begin{equation}\label{eq:curve}
 u(q)=-\frac{2q(1+6q)}{1+8q},\qquad
 v(q)=\frac{16q^4}{1+8q},\qquad B(q)=\HH(u(q),v(q)).
\end{equation}
The target point corresponds to $q_0=1/12$.

\begin{lemma}\label{lem:restriction}
On the curve \eqref{eq:curve},
\begin{equation}\label{eq:gradient}
 4q(2q+1)B+(8q^2+12q+1)U+(24q^2+8q-1)V=0.
\end{equation}
Moreover, $L_BB=0$, where $\tht=\tht_q$ and coefficients in the following
operator are placed to the left of powers of $\tht$:
\begin{align}\label{eq:LB}
L_B={}&-\tht^3
 -2q(22\tht^3+15\tht^2+9\tht+2)\notag\\
 &-4q^2(181\tht^3+219\tht^2+142\tht+32)\notag\\
 &-32q^3(169\tht^3+270\tht^2+181\tht+45)\notag\\
 &-64q^4(244\tht^3+384\tht^2+149\tht-18)\notag\\
 &+512q^5(\tht+1)(20\tht^2+148\tht+159)
 +49152q^6(\tht+1)(\tht+2)(2\tht+3).
\end{align}
The analytic solution of this equation with constant term $1$ is unique.
\end{lemma}

\begin{proof}
We describe an exact certificate for both assertions. The coefficient
ratios of \eqref{eq:H} give $P_1\HH=P_3\HH=0$, with
\begin{align}\label{eq:PDE}
P_1={}&u(1-4u)\partial_u^2+v(1-16u)\partial_u\partial_v
       -12v^2\partial_v^2\notag\\
 &+(1-10u)\partial_u-22v\partial_v-2,\notag\\
P_3={}&u\tht_v^3-v\tht_u(\tht_u+3\tht_v+1)
                              (\tht_u+3\tht_v+2).
\end{align}
For example, if $a_{j,k}$ is the coefficient in \eqref{eq:H}, then
\[
 \frac{a_{j+1,k}}{a_{j,k}}
 =\frac{2(2j+2k+1)(j+3k+1)}{(j+k+1)(j+1)},\qquad
 \frac{a_{j-1,k}}{a_{j,k-1}}
 =\frac{j(j+3k-2)(j+3k-1)}{k^3}.
\]
The second equality is used for $j,k\geq1$; the multipliers in $P_3$
make the boundary terms vanish.

A further, fourth-order radial equation is obtained as follows. Define
\[
 a_n(c)=\sum_{k=0}^n\frac{(n+2k)!}{(n-k)!(k!)^3}c^{n-k},\qquad
 F_c(z)=\sum_{n\geq0}\binom{2n}{n}a_n(c)z^n=\HH(cz,z).
\]
Summing first over $j=n-k$ by the binomial theorem gives
\begin{equation}\label{eq:Gauss}
 A_c(t):=\sum_{n\geq0}a_n(c)t^n
 =\frac1{1-ct}\,{}_2F_1\left(\frac13,\frac23;1;
                              \frac{27t}{(1-ct)^3}\right).
\end{equation}
The pulled-back Gauss equation is
$\sum_{r=0}^4t^rP_r(\tht_t)A_c=0$, where
\begin{align*}
P_0(h)&=-h^2,\\
P_1(h)&=(c+27)h^2+(5c+27)h+c+6,\\
P_2(h)&=3c\bigl((c+18)h^2-2ch-c-2\bigr),\\
P_3(h)&=-c^3h(5h+3),\\
P_4(h)&=2c^4(h+1)^2.
\end{align*}
Here the temporary notation $P_r(h)$ denotes coefficient polynomials,
not the partial differential operators in \eqref{eq:PDE}.
Multiplying coefficients by $\binom{2n}{n}$ gives the fifth-order
annihilator $M=\sum z^rM_r(\tht_z)$, where
\[
 M_r(h)=4^rP_r(h)\frac{(h+\tfrac12)_r}{(h+1)_r}
                         (h+r)(h+r-1)(h+r-2).
\]
Each $M_r(h)$ simplifies to a polynomial. There is an exact Ore
factorization $M=QL$, with
\begin{equation}\label{eq:Ore}
 Q=\frac{-(1+8cz)\tht_z+2(1+2cz)}{1-10cz},\qquad
 L=\sum_{r=0}^4z^r\ell_r(\tht_z),
\end{equation}
where
\begin{align*}
\ell_0(h)&=h^3(h-1),\\
\ell_1(h)&=-2h\bigl((11c+54)h^3+(-c+81)h^2+(7c+39)h+c+6\bigr),\\
\ell_2(h)&=12c(2h+1)\bigl((7c+45)h^3+7ch^2+(6c-5)h+2c\bigr),\\
\ell_3(h)&=-8c^3(2h+1)(2h+3)(17h^2+24h+12),\\
\ell_4(h)&=80c^4(h+1)(2h+1)(2h+3)(2h+5).
\end{align*}
The rule $\tht_z f=f\tht_z+zf'$ verifies \eqref{eq:Ore} coefficient by
coefficient. Thus $E=LF_c$ has the form
$E=Cz^2(1+8cz)^{-3/2}$. The first three coefficients of $F_c$ are
$1,2c+12,6c^2+144c+540$; substitution shows $[z^2]E=0$.
Consequently $LF_c=0$.

Replace $c$ by $u/v$, $z$ by $v$, and $\tht_z$ by
$\tht_u+\tht_v$ in this radial equation. Together with
\eqref{eq:PDE}, it supplies the three equations used by the jet
certificate. More precisely, the certificate gives rational functions
of $q$ expressing $L_BB$ as a linear combination of derivatives of
these three equations of total order at most five. A second multiplier
vector proves, for the left side $W$ of \eqref{eq:gradient},
\begin{equation}\label{eq:Wode}
 \tht_qW=\frac{4q}{1+8q}W.
\end{equation}
Both statements are checked by matching all $21$ ordinary jets
$\partial_u^a\partial_v^b\HH$, $a+b\leq5$, over $\mathbb Q(q)$.
The vectors are given in
\code{prize_congruence_curve90_jet_certificate.json}; the independent
verifier reconstructs the operators and checks every coefficient.
Equation \eqref{eq:Wode} gives $W=C\sqrt{1+8q}$, while $W(0)=0$.
This proves \eqref{eq:gradient}. Finally,
$L_B|_{q=0}=-\tht_q^3$, so the coefficient of $q^n$, $n\geq1$,
is determined uniquely after division by $-n^3$.
\end{proof}

At $q_0$, \eqref{eq:gradient} and the chain rule give
\[
 7B+37U-3V=0,\qquad
 \tht_qB=\frac{14}{15}U+\frac{18}{5}V.
\]
Eliminating $U,V$ therefore reduces \eqref{eq:targetH} to
\begin{equation}\label{eq:reduction}
 S=3\bigl[(35\tht_q+18)B\bigr](1/12).
\end{equation}

\section{The level-60 modular transformation}\label{sec:modparam}
For $\tau$ in the upper half-plane, write $\eta_d=\eta(d\tau)$ and
$Q_\tau=e^{2\pi i\tau}$. Put
$D=(2\pi i)^{-1}\dd/\dd\tau=Q_\tau\dd/\dd Q_\tau$. Define
\begin{align}\label{eq:eta}
 \mathcal A&=\eta_1\eta_3\eta_4\eta_5\eta_{12}\eta_{15}\eta_{20}\eta_{60},&
 \mathcal B&=\eta_2^2\eta_6^2\eta_{10}^2\eta_{30}^2,\notag\\
 X&=\mathcal A/\mathcal B,&
 Z&=\eta_2\eta_6\eta_{10}\eta_{30}/X^2.
\end{align}
Write $\mathcal C=\eta_2\eta_6\eta_{10}\eta_{30}$ and $L=D\log X$.
The product expansions give $X=Q_\tau+O(Q_\tau^2)$,
$Z=1+O(Q_\tau)$, and $L=1+O(Q_\tau)$. The candidate polynomials are
\begin{align}\label{eq:GJ}
 G(x)&=(1-x)(1-4x)(1+4x^2)(1-x+4x^2),\notag\\
 J(x)&=4x(1-8x+40x^2-83x^3+160x^4-144x^5).
\end{align}

\subsection{An independent eta-product and Sturm proof}
\begin{proposition}\label{prop:modparam}
The eta quotients in \eqref{eq:eta} satisfy the identities
\begin{equation}\label{eq:modidentities}
 L^2=G(X)Z^2,\qquad
 2Z D^2Z-3(DZ)^2=Z^4J(X).
\end{equation}
In particular, near the cusp $Q_\tau=0$,
\begin{equation}\label{eq:DXlocal}
 DX=X\sqrt{G(X)}Z,\qquad \sqrt{G(0)}=1.
\end{equation}
As a function of the local coordinate $x=X$, the normalized period
$Z(x)=1+O(x)$ satisfies
\begin{equation}\label{eq:LZ}
 L_ZZ=0,\qquad
 L_Z=G\delta^3+\frac32(\delta G)\delta^2
             +\left(\frac12\delta^2G-J\right)\delta-\frac12\delta J,
 \quad \delta=x\frac{\dd}{\dd x}.
\end{equation}
The analytic solution of \eqref{eq:LZ} with constant term $1$ is unique.
\end{proposition}
\begin{proof}
We give all modularity and holomorphy hypotheses needed for the exact
Sturm checks. For an eta quotient
$f=\prod_{d\mid60}\eta_d^{r_d}$ of integer weight
$k=\tfrac12\sum r_d$, the classical eta-quotient criterion
\cite{Newman,Ligozat} requires
\begin{equation}\label{eq:etacriterion}
 \sum_{d\mid60}d r_d\equiv
 \sum_{d\mid60}(60/d)r_d\equiv0\pmod {24}.
\end{equation}
Its character is the Kronecker character determined by
$(-1)^k\prod d^{r_d}$. Thus, for even $k$, a rational square
$\prod d^{r_d}$ gives trivial character. For
$\mathcal A,\mathcal B,\mathcal C,X,Z$, respectively, the weights are
$4,4,2,0,2$, both sums in \eqref{eq:etacriterion} give the list
$120,96,48,24,0$, and the corresponding products are
$12960000,12960000,3600,1,3600$. These are all squares, and the two
congruences hold for each quotient. Eta has no zeros in the upper
half-plane, so $X$ is a modular unit there.

Put
\[
 S_+=\{1,3,4,5,12,15,20,60\},\qquad S_-=\{2,6,10,30\}.
\]
The exponents of $X$ are $1$ on $S_+$ and $-2$ on $S_-$;
those of $Z$ are $-2$ and $5$, respectively.
The cusp-order formula \cite{Ligozat}, for a cusp of denominator
$c\mid60$, is
\begin{equation}\label{eq:cuspformula}
 v_c(f)=\frac{60}{24\gcd(c^2,60)}
             \sum_{d\mid60}\frac{\gcd(c,d)^2}{d}r_d.
\end{equation}
Here the order is with respect to the local parameter of that cusp.
There is one cusp for each divisor $c$: the multiplicity
$\varphi(\gcd(c,60/c))$ is $1$ in every case. Formula
\eqref{eq:cuspformula} gives the following complete table, grouped
according to the twelve cusp denominators:
\begin{center}
\small
\begin{tabular}{c|rrrr}
denominators $c$&$v_c(X)$&$v_c(Z)$&
$v_c(G(X)Z^2)$&$v_c(J(X)Z^4)$\\ \hline
$2,6,10,30$&$-2$&$6$&$0$&$12$\\
$1,3,4,5,12,15,20,60$&$1$&$0$&$0$&$1$
\end{tabular}
\end{center}
For example, at a pole of $X$, both $G$ and $J$ have degree $6$,
so the last two orders are $6(-2)+2(6)=0$ and
$6(-2)+4(6)=12$. At a zero of $X$, use $G(0)=1$ and the
simple zero of $J$ at $0$. Consequently $Z$ is a holomorphic
weight-$2$ form, and $G(X)Z^2$ and $J(X)Z^4$ are holomorphic forms of
weights $4$ and $8$ on $\Gamma_0(60)$ with trivial character.

The logarithmic derivative of the eta product gives
\begin{equation}\label{eq:LE2}
 L=\frac1{24}\sum_{d\mid60}r_d d E_2(d\tau),
 \qquad
 E_2(\tau)=1-24\sum_{n\geq1}\sigma_1(n)Q_\tau^n,
\end{equation}
where $r_d$ is the exponent vector of $X$ and
$\sigma_1(n)=\sum_{a\mid n}a$. The $E_2$ transformation anomaly
cancels because $\sum r_d=0$, so $L$ has weight $2$.
Alternatively this follows directly by differentiating the weight-$0$
transformation of $X$. It is holomorphic in the upper half-plane.
For any cusp represented by $\sigma\in\mathrm{SL}_2(\mathbb Z)$,
write its width as $w$, and
$X(\sigma\tau)=q_c^m(a_0+O(q_c))$, where
$q_c=e^{2\pi i\tau/w}$ and $a_0\ne0$. Then
\[
 (L|_2\sigma)(\tau)=D\log X(\sigma\tau)
                  =\frac m w+O(q_c).
\]
This proves holomorphy at all cusps, including poles of $X$.
Hence $L^2-G(X)Z^2$ is a holomorphic weight-$4$ form.

For completeness, the derivative expression
$R(f)=2fD^2f-3(Df)^2$ of a holomorphic weight-$2$ form is holomorphic
of weight $8$. It equals one third of the second Rankin--Cohen bracket
$[f,f]_2$. Directly, for
$\gamma=\left(\begin{smallmatrix}a&b\\c&d\end{smallmatrix}\right)$,
$j=c\tau+d$, and $\alpha=c/(2\pi i)$, differentiation gives
\begin{align*}
 f(\gamma\tau)&=j^2f(\tau),\\
 (Df)(\gamma\tau)&=j^4Df(\tau)+2\alpha j^3f(\tau),\\
 (D^2f)(\gamma\tau)&=j^6D^2f(\tau)+6\alpha j^5Df(\tau)
                                  +6\alpha^2j^4f(\tau).
\end{align*}
Substitution cancels the anomaly terms and gives
$R(f)(\gamma\tau)=j^8R(f)(\tau)$. The same calculation for the slash
operator at any cusp gives
$R(f)|_8\sigma=R(f|_2\sigma)$. The right side has a holomorphic local
Fourier series. This proves the cusp holomorphy as well as covariance.
Apply this to $f=Z$. Thus $R(Z)-Z^4J(X)$ is a holomorphic
weight-$8$ form.

The index is
\[
 [\mathrm{SL}_2(\mathbb Z):\Gamma_0(60)]
 =60(1+1/2)(1+1/3)(1+1/5)=144.
\]
The Sturm bounds \cite{Sturm} for the two differences are therefore
$4\cdot144/12=48$ and $8\cdot144/12=96$.
The self-contained program \code{sun90_modparam_verify.py}
proves that their exact integer coefficients vanish from degree $0$
through $48$ and $96$, respectively. It also computes
\eqref{eq:etacriterion}, every cusp order, and the derivative anomaly
cancellation independently from the displayed exponent vectors.

Here is a reproducible specification of the integer-series arithmetic.
For $f=Q_\tau^m\sum_{n\geq0}u_nQ_\tau^n$, $u_0=1$, the eta product
gives
\begin{equation}\label{eq:etarecurrence}
 b_n=-\sum_{\substack{d\mid60\\d\mid n}}r_d d\,\sigma_1(n/d),
 \qquad
 n u_n=\sum_{k=1}^n b_k u_{n-k}.
\end{equation}
Each division is checked exact. Ordinary integer convolutions give the
products, and $D$ multiplies the coefficient of $Q_\tau^n$ by $n$.
The polynomial coefficients used in the checks are
\begin{align*}
 G&=(1,-6,17,-48,68,-96,64),\\
 J&=(0,4,-32,160,-332,640,-576),
\end{align*}
listed in ascending powers. Sturm's theorem now proves both complete
identities in \eqref{eq:modidentities}. Equivalently one can apply its
congruence formulation for every prime to these integral coefficients.
Thus these checks are finite certificates for identities, with their
modularity hypotheses supplied above.

Near the cusp, $L,Z$, and the square root of $G$ all have constant
term $1$, proving \eqref{eq:DXlocal}. Since $X=Q_\tau+O(Q_\tau^2)$,
we may use $x=X$ as a local coordinate. Applying the chain rule to
\eqref{eq:modidentities} gives
\[
 2GZ\delta^2Z-G(\delta Z)^2+(\delta G)Z\delta Z-JZ^2=0.
\]
Its $\delta$ derivative is
\[
 Z\bigl[2G\delta^3Z+3(\delta G)\delta^2Z
             +(\delta^2G-2J)\delta Z-(\delta J)Z\bigr]=0.
\]
Because $Z(0)=1$, division by $2Z$ proves \eqref{eq:LZ}.
Finally its constant-coefficient indicial operator is $\delta^3$:
the coefficient of $x^n$ is determined recursively by division by
$n^3$ for every $n\geq1$. This proves uniqueness.
\end{proof}

\subsection{The algebraic pullback}

Set
\begin{equation}\label{eq:phiK}
 \phi(x)=\frac{x}{2(1-2x)^2},\qquad
 K(x)=(1-2x)^2\sqrt{1+4x^2}.
\end{equation}
An exact differential-operator identity is
\begin{equation}\label{eq:pullback}
 K^{-1}L_B(\phi(x),\tht_q)K
 =\frac{(1+4x^2)^2}{(2x-1)^9(2x+1)}L_Z,
 \qquad \tht_q=\frac{x(1-2x)}{1+2x}\frac{\dd}{\dd x}.
\end{equation}
To verify it, use
$K'/K=-4/(1-2x)+4x/(1+4x^2)$ and compare the four coefficients of
$(\dd/\dd x)^r$, $0\leq r\leq3$. These are rational-function
identities, all checked in the attached verifier. Since
$\phi(0)=0$, $\phi'(0)=1/2$, and $K(0)=1$, uniqueness in
Lemma~\ref{lem:restriction} yields
\begin{equation}\label{eq:BZ}
 B(\phi(x))=K(x)Z(x).
\end{equation}

\section{An independently certified modular equation and CM point}
We next establish the needed special value of $X$, rather than assuming
its entry in a table. Define the symmetric polynomial
\begin{equation}\label{eq:Phi}
 \Phi(x,y)=x^8+y^8+\sum_{i=1}^7\sum_{j=1}^i c_{ij}S_{ij}(x,y),
\end{equation}
where $S_{ii}=x^iy^i$ and $S_{ij}=x^iy^j+x^jy^i$ for $i>j$.
Its complete coefficient array is
\begin{center}
\small
\setlength{\tabcolsep}{5pt}
\begin{tabular}{r|rrrrrrr}
$i\backslash j$&1&2&3&4&5&6&7\\\hline
1&$-1$&&&&&&\\
2&7&$-14$&&&&&\\
3&$-21$&49&$-392$&&&&\\
4&49&$-119$&588&$-168$&&&\\
5&$-77$&231&$-1120$&2352&$-6272$&&\\
6&77&$-245$&924&$-1904$&3136&$-3584$&\\
7&$-56$&308&$-1232$&3136&$-5376$&7168&$-4096$
\end{tabular}
\end{center}

\begin{lemma}\label{lem:CM}
The identity $\Phi(X(\tau),X(7\tau))=0$ holds. At
$\tau_0=i\sqrt{7/60}$,
\begin{equation}\label{eq:CM}
 X(\tau_0)=x_0:=\frac{5-\sqrt{21}}4.
\end{equation}
The branch $y=y(x)$ of $\Phi(x,y)=0$ through $(x_0,x_0)$ satisfies
\begin{equation}\label{eq:ysecond}
 y'(x_0)=-1,\qquad y''(x_0)=\frac{56}{15}+\frac{32\sqrt{21}}{63}.
\end{equation}
\end{lemma}
\begin{proof}
The eta-quotient criterion shows that $\mathcal A$ and $\mathcal B$
are holomorphic weight-$4$ forms of trivial character on $\Gamma_0(60)$.
For their exponent vectors, respectively, both sums
$\sum d r_d$ and $\sum (60/d)r_d$ equal $120$ and $96$; their products
$\prod d^{r_d}$ are squares. All exponents are nonnegative.
Thus
\[
 \mathcal B(\tau)^8\mathcal B(7\tau)^8
                \Phi(X(\tau),X(7\tau))
\]
is a holomorphic weight-$64$ form of trivial character on
$\Gamma_0(420)$. The index of this group is $1152$, so its Sturm
bound is $64\cdot1152/12=6144$. The exact integer eta-product
calculation in \code{prize90_geometry_mod7_sturm.py} proves that
every coefficient from $Q_\tau^0$ through $Q_\tau^{6144}$ is zero.
Sturm's theorem \cite{Sturm} proves the modular identity.
The computation uses exact integer convolution with a rigorous
coefficient bound; its result is not a finite numerical test substituted
for a modular identity.

The eta transformation $\eta(-1/\tau)=\sqrt{-i\tau}\eta(\tau)$,
together with the reciprocal symmetry of the exponents of $X$, gives
\begin{equation}\label{eq:FrickeX}
 X\left(-\frac1{60\tau}\right)=X(\tau).
\end{equation}
Put $\tau_1=i/\sqrt{420}$. Then
$7\tau_1=-1/(60\tau_1)=\tau_0$, so $X(\tau_0)$ lies on the diagonal
of the modular equation. Direct factorization gives
\begin{align}\label{eq:diagonal}
\Phi(x,x)={}&-x^2(2x-1)^2(4x^2-10x+1)(4x^2+2x+1)\notag\\
 &\hspace{10mm}\cdot(4x^3+2x-1)(16x^3-8x^2-1).
\end{align}

For completeness, the correct root is selected by rational inequalities.
Machin's formula with alternating-series bounds, followed by positive
Taylor bounds for the exponential, gives
\[
 \frac{29}{250}<e^{-2\pi\sqrt{7/60}}<\frac{59}{500}.
\]
One may use $6283/2000<\pi<22/7$ and
$683/2000<\sqrt{7/60}<427/1250$, and twenty exponential Taylor terms.
Let $l=29/250$, $h=59/500$, and put
\[
 S_A=\sum_{d\in\{1,3,4,5,12,15,20,60\}}\frac{h^d}{1-h^d},\qquad
 S_B=2\sum_{d\in\{2,6,10,30\}}\frac{h^d}{1-h^d}.
\]
Using $\prod(1-a_n)\geq1-\sum a_n$ first for finite products and then
by a limit, the eta products imply
\[
 l(1-S_A)<X(\tau_0)<\frac{h(1-l)}{1-S_B}.
\]
Exact rational comparison makes the lower bound greater than $1/10$
and the upper bound smaller than $11/100$.
Only one root of \eqref{eq:diagonal} lies in this interval: it is the
smaller root of $4x^2-10x+1$. This proves \eqref{eq:CM}.
These inequalities and the root count are independently checked by
\code{prize90_geometry_cm_verify.py}.

Finally $\Phi_x(x_0,x_0)\ne0$. Symmetry gives $y'(x_0)=-1$, and
implicit differentiation gives
\[
 y''(x_0)=-\frac{2(\Phi_{xx}-\Phi_{xy})}{\Phi_x}(x_0,x_0).
\]
Substitution of the explicit polynomial \eqref{eq:Phi} proves
\eqref{eq:ysecond} in the quadratic field $\mathbb Q(\sqrt{21})$.
\end{proof}

\section{The weighted period value by Fricke differentiation}
Here we derive the weighted evaluation directly, so that no specific
CM table or displayed Ramanujan--Sato formula is needed. Write
\[
 D=\frac1{2\pi i}\frac{\dd}{\dd\tau},\qquad g(x)=\sqrt{G(x)}.
\]
Proposition~\ref{prop:modparam} gives
\begin{equation}\label{eq:DX}
 DX=Xg(X)Z
\end{equation}
on the branch from the cusp. We check its sign at the required CM point
without a table convention. Lemma~\ref{lem:CM} proves
$0<Q_0=e^{-2\pi\sqrt{7/60}}<h=59/500$.
All eta factors, and hence $Z$, are positive for real $0<Q_\tau<1$.
The elementary inequality
\[
 \sum_{n\geq1}\sigma_1(n)t^n\leq\frac{t}{(1-t)^3}
 \quad(0<t<1)
\]
follows from $\sigma_1(n)\leq n(n+1)/2$. Dropping the positive
$S_-$ terms in \eqref{eq:LE2} gives, for $0<Q_\tau\leq h$,
\begin{equation}\label{eq:Lpositive}
 L(Q_\tau)\geq
 1-\sum_{d\in S_+}\frac{d h^d}{(1-h^d)^3}>0.
\end{equation}
The last comparison is exact rational arithmetic, also checked by
\code{sun90_modparam_verify.py}. Thus $L=+\sqrt{G(X)}Z$ at $\tau_0$,
with $g(x_0)>0$. This also shows that $X$ is strictly increasing on the
real nome interval from the cusp to $Q_0$, so that the local coordinate
and period in Proposition~\ref{prop:modparam} reach this CM value on
the chosen branch. The eta transformation
and \eqref{eq:FrickeX} give
\begin{equation}\label{eq:FrickeZ}
 Z\left(-\frac1{60\tau}\right)=-60\tau^2Z(\tau).
\end{equation}
Set $t=\tau_1$, $s=7t=\tau_0$, and abbreviate $Z_s=Z(s)$.
Differentiating \eqref{eq:FrickeX} and \eqref{eq:FrickeZ} yields
\begin{equation}\label{eq:FrickeD}
 Z(t)=7Z_s,\qquad DX(t)=-7DX(s),\qquad
 DZ(t)+49DZ(s)=\frac{14\sqrt{105}}{\pi}Z_s.
\end{equation}
In particular the two occurrences of $DX$ have opposite signs; at $t$
one must use $DX(t)=-x_0g(x_0)Z(t)$.

The branch $y(x)=X(7\tau)$, $x=X(\tau)$, has
\[
 y''(x_0)=\frac{49D^2X(s)+D^2X(t)}{DX(t)^2}.
\]
At $s$ and $t$, respectively, differentiating \eqref{eq:DX} with the
appropriate signs gives
\begin{align*}
 D^2X(s)&=x_0g(g+x_0g')Z_s^2+x_0g\,DZ(s),\\
 D^2X(t)&=x_0g(g+x_0g')Z(t)^2-x_0g\,DZ(t),
\end{align*}
where $g,g'$ are evaluated at $x_0$. Substitute
\eqref{eq:FrickeD}, and use $DZ(s)=x_0gZ_sZ_x(x_0)$. The result is
\begin{equation}\label{eq:yFricke}
 y''(x_0)=\frac2{x_0}+\frac{G'(x_0)}{G(x_0)}
     +\frac{2Z_x(x_0)}{Z_s}
     -\frac{14\sqrt{105}}{49\pi x_0g(x_0)Z_s}.
\end{equation}
Consequently, if
\begin{equation}\label{eq:lambda}
 \lambda=1+\frac{x_0G'(x_0)}{2G(x_0)}-\frac{x_0y''(x_0)}2
         =\frac32-\frac{11\sqrt{21}}{42},
\end{equation}
then \eqref{eq:yFricke} proves the exact evaluation
\begin{equation}\label{eq:value}
 [ (\delta+\lambda)Z ](x_0)
       =\frac{\sqrt{60/7}}{2\pi\sqrt{G(x_0)}}.
\end{equation}
This derivation uses the independent modular equation and its implicit
derivatives from Lemma~\ref{lem:CM}.

\section{Completion of the proof and convergence}
Since $\phi(x_0)=1/12$, differentiate \eqref{eq:BZ} to obtain
\begin{align}\label{eq:weightedbridge}
(35\tht_q+18)B
={}&\frac{35K(x)(1-2x)}{1+2x}\notag\\[-1mm]
 &\cdot\left[\delta+\delta\log K(x)
                   +\frac{18(1+2x)}{35(1-2x)}\right]Z(x).
\end{align}
At $x=x_0$, the constant inside the square brackets equals
$\lambda$ in \eqref{eq:lambda}. All these equalities are exact in
$\mathbb Q(\sqrt{21})$. The positive algebraic factor
\[
 \frac{35K(x_0)(1-2x_0)}{1+2x_0}
               \frac{\sqrt{60/7}}{2\sqrt{G(x_0)}}
\]
has square $900$. Since $0<x_0<1/4$, every displayed square root
and remaining factor is positive, so this factor equals $30$.
Equations \eqref{eq:value} and \eqref{eq:weightedbridge} now give
$[(35\tht_q+18)B](1/12)=30/\pi$. Combined with
\eqref{eq:reduction}, this proves \eqref{eq:main}.

We justify all analytic operations and continuation used above.
Replacing $-324$ by $324$ majorizes the original double series
termwise. From $\binom{2n}{n}\leq4^n$ and \eqref{eq:Gauss}, at
$z_0=1/2160$ the corresponding Gauss parameters are
\[
 1296z_0=\frac35<1,\qquad
 \frac{108z_0}{(1-1296z_0)^3}=\frac{25}{32}<1.
\]
The strict inequalities hold in a neighborhood, proving local uniform
absolute convergence, including differentiated series. A fully rational
majorant is available at $R=1/2088>z_0$:
\[
 1-1296R=\frac{11}{29},\qquad
 \frac{108R}{(1-1296R)^3}=\frac{2523}{2662},\qquad
 \sum_n\left|\binom{2n}{n}a_n(-324)\right|R^n
       \leq\frac{7018}{139}.
\]
For example, the absolute tail of \eqref{eq:main} after $n=N-1$ is at most
\[
 \frac{7018}{139}\left(\frac{29}{30}\right)^N
                 \bigl(30(357N+103)+310590\bigr).
\]
Along $0\leq q\leq1/12$, both $|u(q)|$ and $v(q)$ increase, so this
entire path lies in the same domain of the bivariate series. The leading
coefficient of \eqref{eq:LB} in $\tht_q^3$ is
\[
 (2q-1)(6q+1)(8q+1)^3(16q+1),
\]
which has no zero on $0<q\leq1/12$. For $0\leq x\leq x_0$,
$K(x)$ and $\sqrt{G(x)}$ take their positive branches, and no singularity
of \eqref{eq:LZ} is crossed. The identities proved near the origin
therefore extend to the required point.

\section{Exact certificates and disclosure}
From the repository root, run the complete verification entry point:
\begin{quote}\small\ttfamily
python ramanujan-series/sun90-identity/verify.py
\end{quote}
It stages the programs in \code{code/} and the inputs in
\code{certificates/} into a temporary flat directory, runs all four
components, and leaves tracked files unchanged. Python~3.9 or later and
SymPy~1.14.0 are required. Component JSON reports named below are generated
in that temporary directory; their contents are retained under
\code{component_results} in \code{reproducibility/full_verification.json}.
The ancillary files have distinct roles.
The differential-algebra verification is performed by
\code{verify_sun90.py}, using only
\code{prize_congruence_curve90_jet_certificate.json} alongside it in the
staged directory. This component checks the coefficient-ratio PDEs, the Gauss pullback and Ore
factorization, all jet multipliers, the algebraic modular pullback,
and the final constants and convergence majorant. Its modular special-value
comparison is consistent with \eqref{eq:value}. The general modular data
and the special value are established independently of that comparison
by the eta/Sturm and Fricke proofs above, with the following additional
certificates:
\begin{itemize}
\item \code{sun90_modparam_verify.py} independently computes the eta
      modularity congruences, characters and all twelve cusp orders, checks
      the weight-$4$ and weight-$8$ identities through the complete Sturm
      bounds $48$ and $96$, and verifies the derivative covariance,
      the linear period equation, and the positive sign bound
      \eqref{eq:Lpositive}. It imports no earlier project script.
\item \code{sun90_modparam_result.json} records these exact results,
      the cusp table, and the Fourier-series hashes.
\item \code{sun90_modparam_proof.txt} gives a separate prose account of
      the modularity hypotheses and the same independent proof.
\item \code{prize90_geometry_mod7_result.json} specifies the polynomial
      \eqref{eq:Phi}. The discovery script alone is not its proof.
\item \code{prize90_geometry_mod7_sturm.py} proves the modular identity
      through the complete Sturm bound; its result is recorded in
      \code{prize90_geometry_mod7_sturm_result.json}.
\item \code{prize90_geometry_cm_verify.py} checks the rational nome bounds,
      isolates the CM root, and verifies the implicit derivatives and
      \eqref{eq:lambda}; its result is recorded in
      \code{prize90_geometry_cm_verify_result.json}.
\item \code{verify_sun90_full.py} runs the independent general
      parametrization, algebraic bridge, degree-seven Sturm, and CM
      programs in sequence, checks that their conclusions agree, and
      writes \code{sun90_full_verification.json}.
\end{itemize}
The independent general-parametrization component can also be run alone
from the staged directory with \code{python sun90_modparam_verify.py}. The separate degree-seven Sturm certificate
checks coefficients through degree $6144$,
including the constant term, at weight $64$ and level $420$. Its polynomial
string has SHA-256 digest
\begin{quote}\small\ttfamily
b24469bdc909d58bce0e29b1dc0e51b26\\
a9c47d98d5b19a7f8f6d5cad43fa768.
\end{quote}
The table in \eqref{eq:Phi} also specifies this polynomial without relying
on a digest. No floating-point root recognition is used in the proof.

\paragraph{AI assistance and submission status.}
OpenAI Codex assisted with symbolic exploration, organization of the
argument, and preparation of the exact verification programs and this
manuscript. The author is 李想. The external mathematical
inputs and the executable certificates are disclosed above. The discovery
computations, including exploratory geometric and numerical calculations,
are not assumptions of the proof. Version~1 was emailed to the problem's
author on 16 October 2026. This public Version~2 has not been emailed
separately. Independent external review, first-proof priority, and prize
eligibility remain unconfirmed.

\begin{thebibliography}{9}
\bibitem{Sun}
Z.-W.~Sun, \emph{List of conjectural series for powers of $\pi$ and other
constants}, arXiv:1102.5649v47, 2014.
\url{https://arxiv.org/html/1102.5649}.

\bibitem{vHTY}
M.~van Hoeij, W.-L.~Tsai and D.~Ye,
\emph{The modularity of Z.-W.~Sun's conjectural formulas for $1/\pi$},
arXiv:2410.19289v2, 2024.
\url{https://arxiv.org/html/2410.19289v2}.

\bibitem{CCY}
J.~M.~Campbell, S.~Cooper and D.~Ye,
\emph{Quadratic irrational analogues of Ramanujan's series for $1/\pi$},
arXiv:2602.09352v1, 2026. During discovery, Tables 2--4
suggested the candidate eta quotients and polynomials, and Table~50
provided candidate CM data. The required general identities are
independently proved in Section~\ref{sec:modparam}, and the CM data and
weighted evaluation are independently proved in the following two
sections. No table identity or special value is assumed as a theorem.
\url{https://arxiv.org/html/2602.09352v1}.

\bibitem{Newman}
M.~Newman, \emph{Construction and application of a class of modular
functions (II)}, Proceedings of the London Mathematical Society
\textbf{s3-9} (1959), 373--387.
\url{https://doi.org/10.1112/plms/s3-9.3.373}.

\bibitem{Ligozat}
G.~Ligozat, \emph{Courbes modulaires de genre 1},
Bulletin de la Soci\'et\'e math\'ematique de France, M\'emoire
\textbf{43} (1975), 80~pp.
\url{https://www.numdam.org/item/MSMF_1975__43__5_0/}.

\bibitem{Sturm}
J.~Sturm, \emph{On the congruence of modular forms}, in
\emph{Number Theory (New York, 1984--1985)}, Lecture Notes in Mathematics
\textbf{1240}, Springer, 1987, pp.~275--280.
\end{thebibliography}
\end{document}
